\section{String}
\label{sec:string}

String adalah deretan karakter yang disimpan sebagai byte-byte
berurutan di memori. Setiap karakter menempati satu byte sesuai kode
ASCII. Akhir string dapat ditandai dengan karakter khusus, misalnya
\code{'\$'} untuk layanan DOS fungsi \hex{09}, atau dengan menyimpan
panjangnya secara terpisah.

Bagian ini diringkas dari \cite{gunawan2019}, Bagian 16,
hlm.~154--158, dan \cite{sto2001}, BAB XIX, hlm.~110--116.

\subsection{Deklarasi String}

\begin{lstlisting}[language=MASM, numbers=none]
sumber  DB 'ASSEMBLY$'
tujuan  DB 12 DUP(?)
panjang DB ?
\end{lstlisting}

String \code{'ASSEMBLY\$'} menempati sembilan byte: delapan huruf dan
satu penanda \code{'\$'}. Buffer \code{tujuan} disediakan dua belas
byte agar cukup menampung salinan dan penanda akhir.

\subsection{Menyalin String}

Penyalinan dilakukan byte per byte sampai penanda akhir ditemukan.

\begin{lstlisting}[language=MASM, numbers=none]
LEA SI, sumber
LEA DI, tujuan
XOR CX, CX
SALIN:
    MOV AL, [SI]
    MOV [DI], AL
    CMP AL, '$'
    JE  SELESAI
    INC SI
    INC DI
    INC CX
    JMP SALIN
SELESAI:
    MOV panjang, CL
\end{lstlisting}

\reg{SI} menunjuk sumber dan \reg{DI} menunjuk tujuan. \reg{CX}
menghitung karakter yang disalin, tidak termasuk \code{'\$'}.

\subsection{Membandingkan String}

Perbandingan dilakukan karakter demi karakter. Bila ada perbedaan,
perbandingan dapat dihentikan.

\begin{lstlisting}[language=MASM, numbers=none]
LEA SI, string1
LEA DI, string2
MOV CX, panjang
BANDING:
    MOV AL, [SI]
    CMP AL, [DI]
    JNE BERBEDA
    INC SI
    INC DI
    LOOP BANDING
    ; sampai di sini: string sama
    JMP SELESAI
BERBEDA:
    ; string berbeda
SELESAI:
\end{lstlisting}

\subsection{Mencari Karakter}

Pencarian memakai \asm{CMP} di dalam perulangan. Bila karakter
ditemukan, alamatnya berada pada \reg{SI}.

\begin{lstlisting}[language=MASM, numbers=none]
LEA SI, sumber
CARI:
    MOV AL, [SI]
    CMP AL, '$'
    JE  TIDAK_ADA
    CMP AL, 'M'
    JE  DITEMUKAN
    INC SI
    JMP CARI
DITEMUKAN:
    ; SI menunjuk huruf 'M'
TIDAK_ADA:
\end{lstlisting}

\begin{table}[H]
    \centering
    \caption{Operasi string dan instruksi yang dipakai.}
    \label{tab:operasi-string}
    \small
    \begin{tabular}{@{}p{3.0cm}p{4.6cm}p{5.8cm}@{}}
        \toprule
        \textbf{Operasi} & \textbf{Instruksi utama} &
        \textbf{Keterangan} \\
        \midrule
        Menyalin & \asm{MOV}, \asm{INC}, \asm{CMP}, \asm{JMP} &
        Salin byte sampai penanda akhir \\
        \addlinespace
        Membandingkan & \asm{CMP}, \asm{JNE}, \asm{LOOP} &
        Hentikan bila ada perbedaan \\
        \addlinespace
        Mencari & \asm{CMP}, \asm{JE}, \asm{INC}, \asm{JMP} &
        Hasil berupa alamat pada \reg{SI} \\
        \addlinespace
        Menghitung panjang & \asm{CMP}, \asm{INC}, \asm{JMP} &
        Hitung byte sampai penanda akhir \\
        \bottomrule
    \end{tabular}
\end{table}

\begin{peringatanbox}
\vspace{1.1em}
Buffer tujuan harus cukup besar untuk menampung seluruh karakter
ditambah penanda akhir. Menyalin string tanpa memeriksa batas buffer
dapat menimpa variabel lain di memori.
\end{peringatanbox}

\begin{infobox}
\vspace{1.1em}
String tidak memiliki tipe data khusus. Yang membedakannya dari array
byte biasa hanyalah cara program menafsirkan isinya, misalnya sebagai
kode ASCII dan menggunakan penanda akhir.
\end{infobox}
